% Source: https://artofproblemsolving.com/community/c3h1349080_christmas_math_bee_congrats_jyzhang12
% Stable identity: aops-topic-1349080-q21
The top piece is a cone similar to the whole cone with scale factor $\dfrac{1}{3}$; the top two pieces together form a cone with scale factor $\dfrac{2}{3}$. Volumes of similar solids scale as the cube of the scale factor, so if the top cone has volume $1$ unit, the cones of heights $4$, $8$, $12$ have volumes $1$, $8$, $27$ units.

\begin{asy}
usepackage("amsmath");
size(3.8cm);
defaultpen(fontsize(10pt));
pair V=(0,12);
real[] h={0,4,8};
for(int i=0;i<3;++i){ real y=h[i]; real rr=6*(12-y)/12; draw(ellipse((0,y),rr,0.22*rr),i==0?black:gray); }
draw((-6,0)--V--(6,0),linewidth(0.9));
draw((0,0)--V,dashed);
draw((0,0)--(6,0),red); label("$6$",(3,0),N,red+fontsize(9pt));
draw((8,0)--(8,12)); for(int i=0;i<4;++i) draw((7.7,4i)--(8.3,4i));
for(int i=0;i<3;++i){ draw((6*(12-4i)/12+0.3,4i)--(7.6,4i),dotted); label("$4$",(8,4i+2),E); }
draw((0.3,12)--(7.6,12),dotted);
label("\mbox{}",(min(currentpicture,user=true)+max(currentpicture,user=true))/2,black); // keeps the text after the figure black when LaTeX typesets the labels inline
\end{asy}

The three pieces have volumes $1$, $8-1=7$ and $27-8=19$ units, so the ratio is $\boxed{1:7:19}$.
